\chapter{Literature Survey}
\doublespacing

\section{Introduction}

The literature survey focuses on recent advancements in AI and machine learning for women's healthcare, including menstrual cycle prediction, PCOS detection, and reproductive health monitoring. Various models such as LSTM, XGBoost, and Random Forest have been widely used for accurate prediction and classification. Some studies also explore AI-based chatbots, IoT systems, and telemedicine platforms to improve accessibility and user engagement. Privacy-preserving approaches like federated learning are introduced to protect sensitive health data. However, many existing systems lack integration, real-time monitoring, and personalization. Additionally, issues like small datasets, limited generalization, and lack of explainability remain. These gaps highlight the need for a secure, scalable, and intelligent healthcare solution.

\section{Critical Evaluation of Research Papers}

This section provides a detailed evaluation of all twenty research papers reviewed, examining their contributions, methodologies, and limitations.

\subsection{Paper 1: A Generative Predictive Model for Menstrual Cycle Lengths (2021)}
In this paper, Kathy Li et al. (2021) introduced a generative predictive model for predicting menstrual cycle lengths based on self-tracked mobile health data. The model includes hierarchical structures to handle missing or inconsistent user information and to separate biological signals from tracking artifacts. The model allows for continuous updating and personalization. The proposed method is more accurate than previous ones for predicting menstrual cycle lengths, but it is based on user information and may therefore be biased. \cite{li2021,urteaga2017,li2020}\cite{urteaga2017}\cite{li2020}

\subsection{Paper 2: AI-Based Menstrual Cycle Tracking and Health Assistant (2025)}
An AI-based system for tracking the menstrual cycle using the LSTM algorithm for prediction and Isolation Forest for anomaly detection has been proposed by Anand Kharode et al. in their paper (2025). A chatbot has also been integrated using NLP, which increases the level of engagement and accuracy of the prediction. The proposed system has been developed using the MERN stack, thus ensuring the scalability of the application. However, the privacy of the data has also raised certain concerns. \cite{soni2025}

\subsection{Paper 3: PCOS Disease Classification Using XGBoost (2025)}
In this paper, Atika et al. proposed a machine learning model based on XGBoost, combined with genetic algorithms and SMOTE for classification. The proposed model improves the precision of the classification results and is capable of handling imbalanced data. However, this model may not perform well in different populations. \cite{ahmed2023}

\subsection{Paper 4: Examining Menstrual Tracking to Inform Design (2017)}
Daniel A. Epstein et al. (2017) conducted a study on menstrual tracking using surveys and user feedback. This study provides information on user behavior, user motivation, and the limitations of existing systems. This study highlights the problem of inaccurate predictions and the need for more inclusive designs. This study provides useful insights for designing better applications. This study does not provide any predictive model. \cite{epstein2017,liu2019}

\subsection{Paper 5: Federated AI Framework for Menstrual Health Prediction (2025)}
M. Rajesh (2025) proposed a federated AI system using contactless biosensing for menstrual prediction. This system provides better privacy and personalization using local data processing. This system provides better accuracy and security for real-time predictions. This system uses multimodal data for better performance. This system requires complex infrastructure and hardware. \cite{rajesh2025}

\subsection{Paper 6: ML Methods for Predicting Depression in Menopausal Women (2023)}
Md. Mamun Ali et al. (2023) proposed ML models like Random Forest and XGBoost for depression prediction. This system provides better evaluation of performance metrics and high accuracy. This system helps in the early detection of depression. This study compares different models effectively. This study deals with menopausal women only. \cite{ali2023,isayev2023}

\subsection{Paper 7: Prediction of Fertile Window Using ML (2022)}
In the paper by Jia-Le Yu et al. (2022), the authors used physiological data for prediction. The model has good performance for predicting the fertile window and the menstrual cycles. The model has good performance for individuals with regular cycles. Clinical validation also improves the reliability of the model. However, the performance of the model is not good for individuals with irregular cycles. \cite{yu2022,goodale2019}

\subsection{Paper 8: AI-Based Thermographic Monitoring (2025)}
In the paper by Lucas M. Iturriago-Salas et al. (2025), the authors proposed an AI-driven mobile platform for thermographic maternal health monitoring. The proposed system uses deep learning algorithms on mobile devices such as smartphones. The proposed system has good performance for improving the usability of the system in the healthcare sector. However, the performance of the proposed system is limited by the quality of the image. \cite{iturriago2025}

\subsection{Paper 9: IoT-Based Mobile Health Architecture (2024)}
Ana Gonz\'{a}lez Berm\'{u}dez et al. proposed a scalable mHealth architecture through data fusion and IoT in 2024. The system is useful for developing modular and reusable applications in healthcare. The system is effective for improving interoperability and efficiency. The proposed architecture is applicable to different use cases. However, there is no specific example of implementation. \cite{kilungeja2025}

\subsection{Paper 10: AI Chatbots in Women's Health (2024)}
Hyun-Kyoung Kim presented research on the effects of AI chatbots on women's health in 2024. The study demonstrates positive effects on mental health and users' engagement. Chatbots offer interactive and scalable solutions for users. Chatbots improve the accessibility of healthcare services. However, reliability and accuracy are major concerns. \cite{kim2024}

\subsection{Paper 11: Privacy and Security in Mobile Health Systems (2016)}
David Kotz et al. (2016) studied the security challenges faced by the mobile health system. The paper discusses the risks faced by the system, including data breaches and vulnerabilities. It also emphasizes the significance of protecting the privacy of the users. The handling of the data by the system is of utmost importance, as it can create a sense of trust among the users. However, the paper fails to provide solutions to the problem. \cite{kotz2016}

\subsection{Paper 12: ML Model Comparison for Menstrual Prediction (2025)}
Mutiara Khairunisa et al. (2025) presented a paper that compared the three models, including the LSTM, CNN, and Decision Tree models, for the prediction of the menstrual cycle. The paper shows that the highest accuracy is obtained by the LSTM model, as it can track the patterns of the data. However, the dataset used for the research is small. \cite{khairunisa2025,rego2023}

\subsection{Paper 13: PCOS Prediction Using EHR (2024)}
Zahra Zad et al. (2024) suggested the implementation of machine learning models using EHRs for PCOS prediction. The system uses clinical data such as hormone levels. It provides high prediction accuracy. It can be used for early diagnosis. However, it depends on the availability of EHRs. \cite{zad2024}

\subsection{Paper 14: High-Risk Pregnancy Prediction Using ML (2025)}
Xinyu Pi et al. (2025) implemented machine learning algorithms for predicting high-risk pregnancy. The researchers used multiple classifiers. MLP showed the best results. It provides high accuracy. It can be used for early identification. However, it should be tested using different datasets. \cite{pi2025}

\subsection{Paper 15: Explainable AI in Healthcare (2025)}
Qaiser Abbas et al. (2025) discussed the implementation of explainable AI in clinical decision support systems. It highlights the importance of explainability in AI. It also discusses different techniques used in clinical decision support systems. Explainability helps in making AI more reliable. However, its implementation in real-life scenarios is lacking. \cite{abbas2025}

\subsection{Paper 16: Review of PCOS Detection Techniques (2023)}
The paper presented by Samia Ahmed et al. (2023) is a review-based paper on the detection of PCOS using machine learning algorithms. The paper discusses the performance of algorithms like CNN, SVM, Random Forest, and LR on different datasets. The paper also discusses the limitations of the algorithms, like class imbalance and the unavailability of a standard dataset for PCOS detection. However, the paper does not propose a new model for the detection of PCOS. \cite{ahmed2023}

\subsection{Paper 17: Web-Based ML for PCOS Prediction (2024)}
In the paper titled ``Empowering Early Detection: Web-Based ML Approach for PCOS Prediction'' by Md. Mahbubur Rahman et al. (2024), the author has proposed a web-based approach for the early detection of PCOS using machine learning approaches. The proposed approach has utilized feature selection techniques such as ``Mutual Information'' for better prediction accuracy. The proposed approach has also utilized multiple classifiers such as Random Forest, AdaBoost, and SVM. The proposed approach has also demonstrated high accuracy for the prediction of PCOS. \cite{rahman2024}

\subsection{Paper 18: AI-Based Women's Health Monitoring (2025)}
This paper proposes an AI-driven system for monitoring women's reproductive health and predicting potential disorders using machine learning techniques. The system integrates health parameters such as menstrual patterns, symptoms, and physiological data to provide early risk prediction. It uses supervised learning models to improve prediction accuracy and support decision-making. A key strength is its focus on personalized healthcare and real-time monitoring. However, the system depends on data quality and user input reliability. Overall, it is relevant for developing intelligent healthcare platforms for women's health. \cite{maity2025}

\subsection{Paper 19: AI in Telemedicine (2025)}
Martina Rossi et al. (2025) have presented a systematic review of AI in telemedicine across various domains of healthcare. The paper has mentioned the usage of AI in diagnostics, patient monitoring using wearables, and chatbots in delivering healthcare services. It has also been mentioned that AI can be used to increase the accuracy of diagnostics and to deliver real-time healthcare services. However, there are also challenges that need to be overcome in using AI in telemedicine. The paper has emphasized the importance of AI in delivering telemedicine services. \cite{rossi2025}

\subsection{Paper 20: LLMs in Healthcare (2025)}
Subhankar Maity et al. (2025) have presented a comprehensive review of LLMs in delivering healthcare services. The paper has mentioned the usage of LLMs in delivering clinical decisions, diagnostics, medical education, and patient communication. It has also been mentioned that LLMs can be used to increase the efficiency and knowledge of patients. However, there are also challenges that need to be overcome in using LLMs. The paper has emphasized the importance of LLMs in delivering future AI-driven intelligent healthcare services. \cite{maity2025,wang2024}


\section{Existing Research Paper on Learning}

A significant amount of research has been conducted to enhance privacy and security in federated learning systems. Various techniques such as differential privacy, secure aggregation, and encryption are widely explored to protect sensitive data. Researchers have also analyzed vulnerabilities like data poisoning and inference attacks. These studies aim to improve the reliability and robustness of distributed learning models. The following table summarizes key contributions in this domain.

\begin{center}
\small
\setlength{\tabcolsep}{6pt}
\renewcommand{\arraystretch}{1.2}

\begin{longtable}{|c|p{3cm}|c|p{3.5cm}|p{3cm}|p{3cm}|}
\hline
Sr. & Author & Year & Summary & Advantages & Disadvantages \\
\hline
\endfirsthead

\hline
Sr. & Author & Year & Summary & Advantages & Disadvantages \\
\hline
\endhead

1 & Kathy Li et al. \cite{li2021} & 2021 & Predictive menstrual cycle model & Personalized predictions & Biased data \\
2 & Anand Kharode et al. \cite{soni2025} & 2025 & AI menstrual tracking (LSTM) & Scalable system & Privacy issues \\
3 & Atika et al. \cite{ahmed2023} & 2025 & PCOS classification (XGBoost) & Handles imbalance & Limited generalization \\
4 & Daniel Epstein et al. \cite{epstein2017} & 2017 & Menstrual tracking study & Design insights & No prediction \\
5 & M. Rajesh \cite{rajesh2025} & 2025 & Federated AI health system & High privacy & Complex system \\
6 & Md. Mamun Ali et al. \cite{ali2023} & 2023 & Depression prediction ML & High accuracy & Limited group \\
7 & Jia-Le Yu et al. \cite{yu2022} & 2022 & Fertile window prediction & Good for regular cycles & Poor irregular cycles \\
8 & Lucas Iturriago et al. \cite{iturriago2025} & 2025 & AI thermographic monitoring & Mobile-friendly & Image quality issues \\
9 & Ana Berm\'{u}dez et al. \cite{kilungeja2025} & 2024 & IoT health architecture & Scalable & No implementation \\
10 & Hyun-Kyoung Kim \cite{kim2024} & 2024 & AI chatbots study & Better engagement & Accuracy issues \\
11 & David Kotz et al. \cite{kotz2016} & 2016 & Security in mobile health & Highlights risks & No solutions \\
12 & Mutiara Khairunisa \cite{khairunisa2025} & 2025 & ML comparison (LSTM best) & High accuracy & Small dataset \\
13 & Zahra Zad et al. \cite{zad2024} & 2024 & PCOS using EHR & Early diagnosis & Depends on EHR \\
14 & Xinyu Pi et al. \cite{pi2025} & 2025 & High-risk pregnancy ML & Early detection & Needs validation \\
15 & Qaiser Abbas et al. \cite{abbas2025} & 2025 & Explainable AI & Trust improvement & Limited use \\
16 & Samia Ahmed et al. \cite{ahmed2023} & 2023 & PCOS ML review & Comprehensive & No new model \\
17 & Md. Rahman et al. \cite{rahman2024} & 2024 & Web ML PCOS system & Accurate & Needs testing \\
18 & Sarvesh Gupta et al. \cite{maity2025} & 2025 & AI health monitoring & Personalized & Data dependency \\
19 & Martina Rossi et al. \cite{rossi2025} & 2025 & AI in telemedicine & Better diagnosis & Challenges exist \\
20 & Subhankar Maity et al. \cite{maity2025} & 2025 & LLM healthcare review & Efficient & Ethical issues \\
\hline
\caption{Summary of Existing Research on Learning with Privacy (Papers 1--20)}
\end{longtable}
\end{center}

\section{Critical Evaluation of Existing Research}
\label{sec:critical-eval}

The reviewed literature demonstrates meaningful progress in applying machine learning and AI to women's reproductive health, yet several recurring limitations motivate the development of an integrated platform such as FemCare.

\textbf{Strengths of Existing Approaches.}
Multiple studies confirm that machine learning models are effective for menstrual cycle prediction and PCOS detection. Generative and hierarchical models applied to self-tracked mHealth data have shown state-of-the-art performance for personalised cycle-length prediction \cite{li2021,urteaga2017,li2020}. Physiological and wearable-sensor data have been used to predict fertile windows and detect menstrual phase transitions \cite{yu2022,goodale2019,kilungeja2025}. Gradient-boosted models, including XGBoost, have demonstrated high classification accuracy for PCOS when trained on clinical or EHR datasets \cite{ahmed2023,zad2024,rahman2024}. LLMs and AI chatbots have been shown to produce positive health outcomes in women's health contexts \cite{maity2025,wang2024,kim2024}, while systematic reviews of AI in telemedicine and clinical decision support confirm the broader applicability of these technologies \cite{rossi2025,abbas2025}.

\textbf{Dataset and Population Limitations.}
A consistent limitation across the reviewed studies is the use of small or homogeneous datasets. Several menstrual prediction models were evaluated on datasets that are too limited to generalise across populations with irregular cycles \cite{khairunisa2025,rego2023,li2021}. PCOS prediction models relying on EHR data are restricted to populations that interact with the healthcare system and cannot be applied to community-based screening \cite{zad2024}. High-risk pregnancy prediction models similarly require validation on independent and diverse datasets before clinical application \cite{pi2025}. Studies focused on menopausal or postmenopausal populations \cite{ali2023,isayev2023} further illustrate that findings from one reproductive-age group do not automatically generalise to others.

\textbf{Personalization and Longitudinal Tracking.}
Existing tracking applications demonstrate the importance of self-tracked data for personalised predictions \cite{epstein2017,liu2019}, yet most lack adaptive models that update continuously as new data is recorded. Personalised hormonal-cycle models that separate biological signals from tracking artefacts \cite{li2021,urteaga2017} and approaches that characterise individual physiological variation \cite{li2020,goodale2019} remain largely research prototypes without direct clinical integration.

\textbf{Privacy and Security Gaps.}
Privacy-preserving and federated approaches have been proposed to protect sensitive reproductive health data \cite{rajesh2025,kotz2016}. However, as noted by Kotz et al. \cite{kotz2016}, most mHealth systems still lack robust access control and data isolation mechanisms. Lightweight applications commonly store data locally without database-level per-user isolation, creating vulnerabilities that an integrated cloud-backed system with Row Level Security can address.

\textbf{Explainability Limitations.}
Several studies acknowledge that non-linear ensemble models such as XGBoost and neural networks provide limited interpretability \cite{ahmed2023,zad2024}. Abbas et al. \cite{abbas2025} specifically highlight that explainable AI in clinical settings improves trust but remains difficult to deploy in real-world scenarios. Addressing this gap through feature-importance analysis and SHAP-based explanations is identified as future work for FemCare \cite{lundberg2017}.

\textbf{Integration Gaps.}
Despite strong individual results in cycle prediction \cite{li2021,rego2023,khairunisa2025}, PCOS detection \cite{ahmed2023,zad2024,rahman2024}, chatbot health support \cite{kim2024,maity2025}, and AI-assisted monitoring \cite{iturriago2025,soni2025}, no existing system integrates all of these capabilities alongside healthcare-provider discovery and appointment booking within a single authenticated platform. The FemCare dataset \cite{dataset} was developed specifically to support an integrated prediction and awareness system targeting this gap in women's health technology.

\section{Research Gap}

The available research on women's healthcare is primarily focused on solving individual problems, such as predicting periods or detecting PCOS. \cite{dataset} However, no comprehensive framework is available. Some methods rely on limited data, causing them to be less accurate. Real-time monitoring and personalization for irregular cycles are not yet available. Although methods for ensuring user privacy are available, they have not been effectively implemented. Also, no methods for increasing user trust in AI systems have been proposed yet. Integration of state-of-the-art technologies such as LLMs and telemedicine is yet to be effectively implemented.